\section{Поліноміальна регресія}
Даний вид алгоритму машинного навчання також часто застосовується поряд із
алгоритмом лінійної регресії. У поліноміальній регресії \eqref{eq:formula-5.7}
для побудови залежних змінних регресорів використовуються поліноми $n$-го
степеню. Відповідно чим ступінь такої моделі більший тим краща якість.

\begin{equation}
	\begin{cases}
		b_{0} + b_{1} x + b_{2} x^{2} + \dots + b_{k} x^{k} = y                   \\
		b_{0} x + b_{1} x^{2} + b_{2} x^{3} + \dots + b_{k} x^{k+1}= xy           \\
		b_{0} x^{2} + b_{1} x^{3} + b_{2} x^{4} + \dots + b_{k} x^{k+2}= x^{2} y  \\
		b_{0} x^{3} + b_{1} x^{4} + b_{2} x^{5} + \dots + b_{k} x^{k+3}= x^{3} y  \\
		\dots \dots \dots \dots \dots \dots \dots \dots \dots \dots \dots \dots   \\
		b_{0} x^{k} + b_{1} x^{k+1}+ b_{2} x^{k+2}+ \dots + b_{k} x^{2k}= x^{k} y
	\end{cases}
	\label{eq:formula-5.7}
\end{equation}

Як ми бачимо для кожної можливої моделі будується поліноміальна залежність між
параметрами b та x. Таким чином перед подачею на вхід лінійної регресії подається
не звичайні регресори, а вже у вигляді поліномів.

До переваг використання поліноміальних моделей можна віднести те, що поліноми забезпечують
найкраще наближення залежності між залежною та незалежною змінною,може приймати великий
спектр функцій і взагалі поліном в основному підходить для широкого діапазону кривизни.

Щодо недоліків даного алгоритму є можлива наявність асистентних даних, що в при
роботі методу може серйозно вплинути на результати нелінійного аналізу. Також даний
тип регресійних моделей занадто чутливі до лишків. Крім того, на жаль, є менше
інструментів валідації моделей для виявлення відшаровувачів у нелінійній
регресії, ніж для лінійної регресії.
